\documentclass[11pt]{article}
% ======================================================================
%  weekpreamble.tex — shared preamble for the weekly course summaries (EN)
%  Concise, readable, intuition-first. Same box style as the main doc.
%  Compiles with pdflatex.
% ======================================================================
\usepackage[utf8]{inputenc}
\usepackage[T1]{fontenc}
\usepackage[english]{babel}
\usepackage[a4paper,margin=2.2cm]{geometry}
\usepackage{amsmath,amssymb,amsthm,mathtools}
\usepackage{bm}
\usepackage{enumitem}
\usepackage{xcolor}
\usepackage[most]{tcolorbox}
\usepackage{microtype}

\emergencystretch=3em
\allowdisplaybreaks[1]
\setlength{\parindent}{0pt}
\setlength{\parskip}{0.45em}

% ---------- math macros (same names as the main document) ----------
\newcommand{\E}{\mathbb{E}}
\DeclareMathOperator{\Var}{Var}
\DeclareMathOperator{\Cov}{Cov}
\DeclareMathOperator{\Corr}{Corr}
\DeclareMathOperator*{\argmin}{arg\,min}
\DeclareMathOperator{\MSE}{MSE}
\DeclareMathOperator{\trace}{tr}
\newcommand{\Reals}{\mathbb{R}}
\newcommand{\Ints}{\mathbb{Z}}
\newcommand{\Comp}{\mathbb{C}}
\newcommand{\Nat}{\mathbb{N}}
\newcommand{\Prob}{\mathbb{P}}
\newcommand{\Tr}{^{\mathsf{T}}}
\newcommand{\Herm}{^{\mathsf{H}}}
\newcommand{\conj}[1]{\overline{#1}}
\newcommand{\abs}[1]{\left\lvert #1 \right\rvert}
\newcommand{\norm}[1]{\left\lVert #1 \right\rVert}
\newcommand{\ind}{\mathbf{1}}
\newcommand{\eps}{\varepsilon}
\newcommand{\diff}{\nabla}
\newcommand{\acvs}{\gamma}
\newcommand{\acf}{\rho}
\newcommand{\pacf}{\alpha}
\newcommand{\sdf}{\mathcal{S}}
\newcommand{\filt}{\mathcal{L}}
\newcommand{\Bsh}{B}
\newcommand{\ms}{\;\overset{\mathrm{ms}}{=}\;}
\newcommand{\toprob}{\xrightarrow{P}}
\newcommand{\todist}{\xrightarrow{d}}
\newcommand{\iid}{\overset{\text{iid}}{\sim}}

\setlist[itemize]{leftmargin=1.5em,topsep=2pt,itemsep=2pt}
\setlist[enumerate]{leftmargin=1.7em,topsep=2pt,itemsep=2pt}

% ---------- compact, titled (unnumbered) boxes ----------
\tcbset{wkbase/.style={enhanced,breakable,fonttitle=\bfseries,coltitle=white,
  boxrule=0.6pt,arc=1.2mm,left=2.2mm,right=2.2mm,top=1.1mm,bottom=1.1mm,
  before skip=6pt,after skip=6pt,
  attach boxed title to top left={yshift=-3.2mm,xshift=2.4mm},
  boxed title style={boxrule=0pt,arc=0.5mm}}}

\newtcolorbox{defn}[1]{wkbase, colback=blue!4!white, colframe=blue!58!black,
  colbacktitle=blue!58!black, title={Definition --- #1}, top=3mm}
\newtcolorbox{result}[1]{wkbase, colback=red!4!white, colframe=red!60!black,
  colbacktitle=red!60!black, title={#1}, top=3mm}
\newtcolorbox{intu}[1][Intuition]{wkbase, colback=yellow!12!white,
  colframe=orange!72!black, colbacktitle=orange!72!black, coltitle=white,
  title={#1}, top=3mm}
\newtcolorbox{retenir}[1][Key takeaways --- the week's reflexes]{wkbase,
  colback=green!5!white, colframe=green!48!black, colbacktitle=green!48!black,
  title={#1}, top=3mm}
\newtcolorbox{exmpl}[1][Worked example]{wkbase, colback=black!3!white,
  colframe=black!55!white, colbacktitle=black!55!white, title={#1}, top=3mm}

% ---------- week header ----------
\newcommand{\weekheader}[4]{%
  {\small\sffamily\bfseries\color{black!60} TIME SERIES~~$\cdot$~~MATH-342, EPFL}\par
  \vspace{2pt}
  {\LARGE\bfseries Week #1 --- #3\par}
  \vspace{1pt}
  {\itshape\color{black!55} #2}\par
  \vspace{6pt}
  \begin{tcolorbox}[enhanced,colback=blue!3!white,colframe=blue!45!black,
    arc=1.4mm,boxrule=0.7pt,left=3mm,right=3mm,top=2mm,bottom=2mm,before skip=2pt,after skip=10pt]
    \textbf{The big picture.}~ #4
  \end{tcolorbox}
}

\begin{document}
\weekheader{12}{14 May 2026}{Non-Linear Time Series: ARCH \& GARCH}{For the changing volatility of financial returns (volatility clustering) we model the conditional variance: ARCH/GARCH are uncorrelated yet dependent.}

Linear models (AR, MA, ARMA) describe the conditional \emph{mean} and assume a constant innovation
variance. Financial log-returns break both habits: they are nearly \emph{uncorrelated} (a
conditional-mean model is almost useless), yet their \emph{squares} are strongly dependent --- big
moves cluster with big moves, calm with calm. This \textbf{volatility clustering} plus heavy
(\textbf{fat}) tails calls for a model of the \emph{conditional variance}. ARCH/GARCH keep the
series mean-zero white noise while letting $\sigma_t^2$ evolve with past squared returns (and, for
GARCH, past variances). Throughout, $\mathcal F_t$ is the information up to time $t$, and
$\{\eps_t\}$ is white noise with $\E[\eps_t]=0$, $\Var(\eps_t)=1$, independent of $\mathcal F_{t-1}$.

\section*{Key definitions}

\begin{defn}{Log-returns \& volatility clustering}
For prices $\{X_t\}$, the \textbf{log-returns} are $r_t=\log(X_t/X_{t-1})=\log X_t-\log X_{t-1}$
(relative price changes). Empirically: (i) the ACF/PACF of $r_t$ looks like \emph{white noise};
(ii) the ACF/PACF of $r_t^2$ \emph{decays slowly} $\Rightarrow$ the conditional variance changes
over time; (iii) \textbf{volatility clustering} (large changes follow large, small follow small);
(iv) \textbf{fat tails}, kurtosis $>3$ (a QQ-plot against $\mathcal N(0,1)$ bends at the ends).
\end{defn}

\begin{defn}{ARCH$(p)$ model}
$\{r_t\}$ is \textbf{AutoRegressive Conditionally Heteroscedastic} of order $p$ if
\[
r_t=\sigma_t\eps_t,\qquad
\sigma_t^2=\alpha_0+\sum_{j=1}^{p}\alpha_j\,r_{t-j}^2,
\]
with $\alpha_0>0$, $\alpha_p>0$, $\alpha_j\ge0$. The volatility $\sigma_t$ is
$\mathcal F_{t-1}$-measurable (known one step ahead), so
$\E[r_t\mid\mathcal F_{t-1}]=0$ and $\Var(r_t\mid\mathcal F_{t-1})=\sigma_t^2$. The case $p=1$ is
$\sigma_t^2=\alpha_0+\alpha_1 r_{t-1}^2$.
\end{defn}

\begin{defn}{GARCH$(p,q)$ model}
$\{r_t\}$ is \textbf{Generalised ARCH} if $r_t=\sigma_t\eps_t$ with
\[
\sigma_t^2=\alpha_0+\sum_{j=1}^{p}\alpha_j\,r_{t-j}^2+\sum_{j=1}^{q}\beta_j\,\sigma_{t-j}^2,
\]
$\alpha_0>0$, $\alpha_p,\beta_q>0$, other $\alpha_j,\beta_j\ge0$. Feeding back past
\emph{variances} $\sigma_{t-j}^2$ makes volatility persistent far more cheaply than a large-$p$
ARCH. A GARCH$(p,0)$ is exactly an ARCH$(p)$; again $\Var(r_t\mid\mathcal F_{t-1})=\sigma_t^2$.
\end{defn}

\begin{defn}{Standardised residuals (diagnostics)}
After fitting, $\tilde e_t=r_t/\hat\sigma_t$. If the model is right these look i.i.d.\ $(0,1)$:
the ACF/PACF of $\tilde e_t$ \emph{and} of the squares $\tilde e_t^2$ should resemble white noise,
and a QQ-plot against the assumed innovation law should be roughly straight. One also overlays the
in-sample bands $\pm1.96\,\hat\sigma_t$ on $r_t$ (Gaussian innovations).
\end{defn}

\begin{defn}{ARMA--GARCH model}
Since a GARCH process is itself white noise, it can be the innovation of an ARMA: with $r_t$ a
GARCH process,
\[
X_t=\sum_{j=1}^{p}\phi_j X_{t-j}-\sum_{k=0}^{q}\theta_k\,r_{t-k},\qquad \theta_0=-1,
\]
the $\phi_j,\theta_k$ satisfying the usual ARMA conditions. This models a conditional mean (ARMA)
and a time-varying conditional variance (GARCH) at once --- just be careful with the likelihood, as
errors are no longer i.i.d.\ Gaussian.
\end{defn}

\begin{intu}[Why feed back $\sigma_{t-1}^2$?]
In ARCH(1) the variance depends only on the \emph{single} most recent $r_{t-1}^2$: one quiet day
collapses $\sigma_t^2$ back toward $\alpha_0$, so the model cannot \emph{sustain} a high-volatility
regime. GARCH adds the smooth term $\beta_1\sigma_{t-1}^2$: yesterday's variance estimate is mostly
carried forward, so volatility is \textbf{sticky}. Fitted GARCH$(1,1)$ typically has
$\alpha_1+\beta_1$ close to $1$ --- long, persistent memory (still geometric decay, but very slow).
\end{intu}

\section*{Key results \& formulas}

\begin{result}{Proposition --- moments of ARCH(1)}
For ARCH(1) with $\alpha_1<1$:
\[
\begin{aligned}
\E[r_t\mid\mathcal F_{t-1}]&=0, & \Var(r_t\mid\mathcal F_{t-1})&=\sigma_t^2=\alpha_0+\alpha_1 r_{t-1}^2,\\
\E[r_t]&=0, & \Var(r_t)&=\frac{\alpha_0}{1-\alpha_1}.
\end{aligned}
\]
\textit{Idea:} $\sigma_t$ is known given $\mathcal F_{t-1}$ and $\eps_t$ is independent of it, so
$\E[r_t\mid\mathcal F_{t-1}]=\sigma_t\E[\eps_t]=0$ and $\E[r_t^2\mid\mathcal F_{t-1}]=\sigma_t^2$;
take unconditional expectations and use stationarity $\Var(r_t)=\Var(r_{t-1})$. Here $\alpha_1<1$ is
\emph{necessary and (for ARCH(1)) sufficient} for a finite, positive variance.
\end{result}

\begin{result}{Proposition --- (G)ARCH is white noise (uncorrelated but dependent)}
Any ARCH$(p)$ or GARCH$(p,q)$ is mean-zero and serially uncorrelated: for $\tau\ne0$,
\[
\E[r_t]=0,\qquad \Cov(r_{t+\tau},r_t)=0,\qquad \Corr(r_{t+\tau},r_t)=0 .
\]
\textit{Idea:} for $\tau>0$, both $r_t$ and $\sigma_{t+\tau}$ are $\mathcal F_{t+\tau-1}$-measurable,
so $\E[r_{t+\tau}r_t]=\E\big[r_t\sigma_{t+\tau}\,\E[\eps_{t+\tau}\mid\mathcal F_{t+\tau-1}]\big]=0$.
The qualitative punchline: (G)ARCH leaves the \emph{linear} correlation flat (white noise) yet
creates strong \emph{nonlinear} dependence in the squares. Mean and variance modelling decouple.
\end{result}

\begin{result}{Theorem --- Isserlis (Wick formula) \& ARCH(1) kurtosis}
For a zero-mean Gaussian vector, $\E[X_1X_2X_3X_4]=\E[X_1X_2]\E[X_3X_4]+\E[X_1X_3]\E[X_2X_4]+\E[X_1X_4]\E[X_2X_3]$; in particular $\E[\eps_t^4]=3$. Hence, for ARCH(1) with Gaussian noise, the
fourth moment is finite \textbf{iff} $\alpha_1^2<\tfrac13$, and
\[
\E[r_t^4]=\frac{3\alpha_0^2(1+\alpha_1)}{(1-\alpha_1)(1-3\alpha_1^2)},\qquad
\kappa=\frac{\E[r_t^4]}{\Var(r_t)^2}=\frac{3(1-\alpha_1^2)}{1-3\alpha_1^2}>3\ \ (\alpha_1>0).
\]
So ARCH \emph{generates} excess kurtosis (fat tails) on its own, even with Gaussian $\eps_t$.
\end{result}

\begin{result}{Proposition --- ACF of the squares of ARCH(1)}
For ARCH(1) with $0<\alpha_1<1/\sqrt3$,
\[
\Corr(r_{t+\tau}^2,r_t^2)=\alpha_1^{\abs\tau},\qquad \tau\in\Ints :
\]
the autocorrelation of the \emph{squares} decays \textbf{geometrically}.
\textit{Idea:} condition on $\mathcal F_{t+\tau-1}$ to peel one lag, giving the recursion
$\Cov(r_{t+\tau}^2,r_t^2)=\alpha_1\Cov(r_{t+\tau-1}^2,r_t^2)$ (constant terms cancel via
$\E[r_t^2]=\alpha_0/(1-\alpha_1)$). This is often \emph{too fast} for real data --- motivating GARCH.
\end{result}

\begin{result}{Proposition --- GARCH$(p,q)$ stationarity, variance \& 4th moment}
A GARCH$(p,q)$ has a weakly stationary solution \textbf{iff} $\sum_j\alpha_j+\sum_j\beta_j<1$, and
then
\[
\Var(r_t)=\frac{\alpha_0}{1-\sum_{j=1}^{p}\alpha_j-\sum_{j=1}^{q}\beta_j}.
\]
For GARCH$(1,1)$ with Gaussian noise the fourth moment exists \textbf{iff}
$3\alpha_1^2+2\alpha_1\beta_1+\beta_1^2<1$ (the analogue of $\alpha_1^2<1/3$). \emph{Note:}
$\sum\alpha+\sum\beta<1$ is also sufficient for strict stationarity; the boundary
$\sum\alpha+\sum\beta=1$ (Gaussian noise) is strictly --- not weakly --- stationary: the IGARCH case.
\end{result}

\begin{result}{Proposition --- conditional-Gaussian MLE}
Conditioning on $r_1$ (so no marginal law for the initial value), the likelihood factorises and
each $r_t\mid\mathcal F_{t-1}\sim\mathcal N(0,\sigma_t^2)$. Up to constants,
\[
\log L_c=-\tfrac12\sum_{t}\Big[\log\sigma_t^2+\frac{r_t^2}{\sigma_t^2}\Big],
\]
maximised numerically (with $\sigma_t^2$ from the model recursion) subject to positivity and
$\sum\alpha+\sum\beta<1$. The per-time variance $\sigma_t^2$ is what is fitted, not a constant.
\end{result}

\section*{Intuition}

\begin{intu}[White noise $\ne$ no structure]
The headline paradox: $\Corr(r_{t+\tau},r_t)=0$ (uncorrelated, so linear forecasting of $r_t$ is
hopeless) \emph{but} $\Corr(r_{t+\tau}^2,r_t^2)\ne0$ (the magnitudes are predictable). ARCH/GARCH
live exactly here --- flat ACF, dependent squares. That is also why they slot in cleanly as ARMA
innovations: they do not disturb the correlation structure.
\end{intu}

\begin{intu}[Two nested thresholds --- don't confuse them]
$\alpha_1<1$ buys a finite \emph{variance}; the strictly tighter $\alpha_1^2<\tfrac13$ (i.e.\
$\alpha_1<1/\sqrt3\approx0.577$) buys a finite \emph{fourth} moment, needed before any kurtosis or
ACF-of-squares formula is valid. As $\alpha_1\uparrow1/\sqrt3$, $\kappa\to\infty$: more ARCH effect
$\Rightarrow$ heavier tails. (GARCH$(1,1)$ analogue: $3\alpha_1^2+2\alpha_1\beta_1+\beta_1^2<1$.)
\end{intu}

\begin{intu}[GARCH(1,1) $\approx$ ARCH$(\infty)$, but parsimonious]
Iterating $\sigma_t^2=\alpha_0+\alpha_1 r_{t-1}^2+\beta_1\sigma_{t-1}^2$ gives
$\sigma_t^2=\tfrac{\alpha_0}{1-\beta_1}+\alpha_1\sum_{k\ge0}\beta_1^k r_{t-1-k}^2$ --- an
ARCH$(\infty)$ with geometrically decaying weights $\alpha_1\beta_1^k$. Three numbers
$(\alpha_0,\alpha_1,\beta_1)$ encode the whole slow tail: $\alpha_1$ = immediate shock impact,
$\beta_1$ = persistence/decay rate. An ARCH$(p)$ would need many $\alpha_j$ to mimic this.
\end{intu}

\begin{intu}[Reading the (G)ARCH likelihood]
Each factor is a normal with its \emph{own} variance $\sigma_t^2$. The log-likelihood trades off
``don't make $\sigma_t^2$ needlessly large'' (the $\log\sigma_t^2$ penalty) against ``explain a big
$r_t^2$ with a big $\sigma_t^2$'' (the $r_t^2/\sigma_t^2$ term). That balancing act \emph{is}
volatility tracking.
\end{intu}

\begin{intu}[Pros \& cons]
\textbf{Pros:} capture volatility clusters and time-varying variance; reproduce fat tails (large
kurtosis); GARCH gives persistent volatility cheaply; they don't disturb the white-noise
correlation. \textbf{Cons:} positive and negative shocks act \emph{symmetrically} (no leverage
effect); volatility dependence decays only \emph{geometrically}, often too fast for real markets.
Extensions (EGARCH, GJR-GARCH, long-memory GARCH) fix these but are out of scope.
\end{intu}

\begin{exmpl}[GARCH(1,1) sanity checks (S\&P 500 fit)]
Fitted: $\alpha_0=3.686\times10^{-6}$, $\alpha_1=0.1713$, $\beta_1=0.7897$.
\emph{Stationary?} $\alpha_1+\beta_1=0.961<1$ \checkmark, with
$\Var(r_t)=\alpha_0/(1-0.961)\approx9.45\times10^{-5}$, i.e.\ daily volatility
$\sqrt{\cdot}\approx0.97\%$. \emph{Fourth moment?}
$3\alpha_1^2+2\alpha_1\beta_1+\beta_1^2\approx0.982<1$ \checkmark (only just --- very heavy tails).
The persistence $\alpha_1+\beta_1=0.961$ near $1$ is the hallmark of financial GARCH: long
quiet/turbulent regimes, yet still stationary.
\end{exmpl}

\section*{Key takeaways}

\begin{retenir}
\begin{itemize}
  \item \textbf{Model the variance, not the mean.} $r_t=\sigma_t\eps_t$; ARCH:
        $\sigma_t^2=\alpha_0+\sum\alpha_j r_{t-j}^2$; GARCH adds $\sum\beta_j\sigma_{t-j}^2$.
        Always $\Var(r_t\mid\mathcal F_{t-1})=\sigma_t^2$.
  \item \textbf{White noise but dependent:} $\Corr(r_{t+\tau},r_t)=0$ yet
        $\Corr(r_{t+\tau}^2,r_t^2)=\alpha_1^{\abs\tau}$ (ARCH(1)) $\ne0$. Diagnose with the ACF of
        \emph{squared} residuals.
  \item \textbf{Long-run variance / stationarity in one move:} take unconditional expectations of the
        $\sigma_t^2$ recursion, set $\E[r_{t-j}^2]=\E[\sigma_{t-j}^2]=\Var(r_t)$, solve
        $\Var(r_t)=\alpha_0/(1-\sum\alpha-\sum\beta)$; positivity $\Leftrightarrow\sum\alpha+\sum\beta<1$.
  \item \textbf{Two thresholds:} $\alpha_1<1$ (variance) vs.\ $\alpha_1^2<1/3$ (kurtosis); GARCH(1,1)
        4th moment: $3\alpha_1^2+2\alpha_1\beta_1+\beta_1^2<1$.
  \item \textbf{Fat tails for free:} Gaussian $\eps_t$ + ARCH $\Rightarrow\kappa>3$.
  \item \textbf{Fit by MLE:} maximise $-\tfrac12\sum_t[\log\sigma_t^2+r_t^2/\sigma_t^2]$ with
        $\sigma_t^2$ from the recursion. GARCH plugs into ARMA as the innovation (ARMA--GARCH).
\end{itemize}
\end{retenir}

\end{document}
